%------------------------------------------------------------------------------%

\usepackage{calculo_varias_variaveis-1}

\title{Multiplicadores de Lagrange}

%------------------------------------------------------------------------------%
\begin{document}

\addtitle

%------------------------------------------------------------------------------%
\section{Extremos Condicionados}

%------------------------------------------------------------------------------%
\begin{frame}{Extremos Condicionados}
  Encontre o máximo (mínimo) da função \(f(x, y, z)\)\\[2mm]
  sujeito a \(g(x, y, z) = 0\)

  \pause
  \vspace{2\baselineskip}
  \emph{$f$} função objetivo

  \pause
  \emph{$g$} restrição
\end{frame}

%------------------------------------------------------------------------------%
%\framedgraphic{Extremos Condicionados}{campo_gradiente}{}

%------------------------------------------------------------------------------%
\section{Extremos Locais Restritos a uma Curva}

%------------------------------------------------------------------------------%
\begin{frame}{Teorema do Gradiente Ortogonal}
  Seja \(f(x, y, z)\) uma função diferenciável em uma região

  \pause
  cujo interior contém a curva
  \[
    C\colon r(t)
    = g(t)\veci
    + h(t)\vecj
    + k(t)\veck
  \]

  \pause
  Se $P_0$ é um ponto em $C$ onde $f$ possua um máximo ou mínimo local relativo
  aos seus valores em $C$,
  \pause
  então \(\nabla f\) é ortogonal a $C$ em $P_0$
\end{frame}

%------------------------------------------------------------------------------%
\begin{frame}{Demonstração}
  Os valores de $f$ em $C$ são dados pela composta
  \[
    p(t) = f(g(t), h(t), k(t))
  \]
  \pause
  então
  \[
    \frac{dp}{dt}
    \pause
    = \D{f}{g}\frac{dg}{dt}
    + \D{f}{h}\frac{dh}{dt}
    + \D{f}{k}\frac{dk}{dt}
    \pause
    = \nabla f \cdot \frac{dr}{dt}
  \]

  \pause
  Se $f$ possui máximo (mínimo) em \((a, b, c)\) correspondente a $t_0$
  \[
    \nabla f(a, b, c) \cdot \frac{dr}{dt}(t_0) = 0
  \]
\end{frame}

%------------------------------------------------------------------------------%
\section{Multiplicadores de Lagrange}

\framedgraphic{Extremos Condicionados}{LagrangeMultipliers2D}{}

%------------------------------------------------------------------------------%
\begin{frame}{Método dos Multiplicadores de Lagrange}
  Suponha que \(f(x, y, z)\) e \(g(x, y, z)\) sejam diferenciáveis\\[2mm]
  e \(\nabla g\neq 0\) quando \(g(x, y, z)=0\)

  \pause
  \vspace{\baselineskip}
  Para encontrar os valores extremos locais de \(f(x, y, z)\) \\[2mm]
  \pause
  sujeito a restrição \(g(x, y, z)=0\) \\[2mm]
  \pause
  encontre $x$, $y$, $z$ e $\lambda$ tais que
  \[
    \nabla f = \lambda \nabla g
    \qquad\qquad
    g(x, y, z) = 0
  \]
  \pause
  \(\lambda\) é o Multiplicador de Lagrange
\end{frame}

%------------------------------------------------------------------------------%
%------------------------------------------------------------------------------%
\begin{question}
  Encontre o maior e menor valores que a função
  \[
    f(x, y) = xy
  \]
  assume na elipse
  \[
    \frac{x^2}{8} + \frac{y^2}{2} = 1
  \]
\end{question}

%------------------------------------------------------------------------------%
\begin{resolution}{Interpretando o Problema}
  Queremos encontrar os valores extremos da função
  \[
    f(x, y) = xy
  \]
  \pause
  sujeito a restrição
  \[
    g(x, y) = \frac{x^2}{8} + \frac{y^2}{2} - 1 = 0
  \]

  \pause
  Para isso precisamos encontrar os valores $x$, $y$ e $\lambda$ tais que
  \[
    \nabla f = \lambda \nabla g
    \qquad\qquad
    g(x, y) = 0
  \]
\end{resolution}

%------------------------------------------------------------------------------%
\begin{resolution}{Gradientes}
\begin{columns}
\column{0.4\linewidth}
  Gradiente de $f$
  \[
    \pause
    \D{f}{x}
    \pause
    = \D{}{x}\left(xy\right)
    \pause
    = y
  \]
  \[
    \pause
    \D{f}{y}
    \pause
    = \D{}{y}\left(xy\right)
    \pause
    = x
  \]
  \pause
  \[
    {\color<19>{blue}\nabla f = \vetor{y}{x}}
  \]
\column{0.6\linewidth}
  \pause
  Gradiente de $g$
  \[
    \pause
    \D{g}{x}
    \pause
    = \D{}{x}\left(\frac{x^2}{8} + \frac{y^2}{2} - 1\right)
    \pause
    = \frac{2x}{8}
    \pause
    = \frac{x}{4}
  \]
  \[
    \pause
    \D{g}{y}
    \pause
    = \D{}{y}\left(\frac{x^2}{8} + \frac{y^2}{2} - 1 \right)
    \pause
    = \frac{2y}{2}
    \pause
    = y
  \]
  \pause
  \[
    {\color<19>{blue}\nabla g = \vetor{\dfrac{x}{4}}{y}}
  \]
\end{columns}
\end{resolution}

%------------------------------------------------------------------------------%
\begin{resolution}{Sistema}
\begin{columns}[t]
\column{0.4\linewidth}
  A equação
  \[
    \nabla f = \lambda \nabla g
  \]
  \pause
  se torna
  \[
    \vetor{y}{x} = \lambda \Vetor{\dfrac{x}{4}}{y}
  \]
  {\color<8>{blue}%
  \pause
  \[
    \begin{cases}
      4y = \lambda x \\[2mm]
      x = \lambda y
    \end{cases}
  \]}
\column{0.4\linewidth}
  \pause
  A equação
  \[
    g(x, y) = 0
  \]
  \pause
  se torna
  \[
    \frac{x^2}{8} + \frac{y^2}{2} - 1 = 0
  \]
  \[
    \pause
    \frac{x^2}{8} + \frac{y^2}{2} = 1
  \]
  {\color<8>{blue}%
  \pause
  \[
    x^2 + 4y^2 = 8
  \]}
\end{columns}
\end{resolution}

%------------------------------------------------------------------------------%
\begin{resolution}{Resolvendo o Sistema Não-Linear}
\begin{columns}
\column{0.5\linewidth}
  Precisamos resolver o sistema
  \[
    \begin{cases}
      4y = \lambda x \\
      x  = \lambda y \\
      x^2 + 4y^2 = 8
    \end{cases}
  \]
  \pause
  Substituindo \(x = \lambda y\) em \( 4y = \lambda x\)
  \[
    \pause
    4y = \lambda x
    \pause
    = \lambda \lambda y
  \]
  \[
    \pause
    4y = \lambda^2 y
    \pause
    \qquad
    \text{\color{red}Não divida por $y$!}
  \]
\column{0.5\linewidth}
  \pause
  Assim \(y=0\)
  \pause
  ou \(y\neq 0\)
  \pause
  e
  \begin{align*}
    \onslide<+->{\lambda^2     & = 4  \\}
    \onslide<+->{\abs{\lambda} & = 2  \\}
    \onslide<+->{\lambda       & = \pm 2}
  \end{align*}

  \onslide<+->{Caso 1: \emph{\(y=0\)}} \\
  \onslide<+->{Caso 2: \emph{\(y\neq 0\)} e \emph{\(\lambda=2\)}} \\
  \onslide<+->{Caso 3: \emph{\(y\neq 0\)} e \emph{\(\lambda=-2\)}}
\end{columns}
\end{resolution}

%------------------------------------------------------------------------------%
\begin{resolution}{Caso 1}
  \[
    y = 0
  \]
  \pause
  \[
    x = \lambda y
    \pause
    = 0
  \]
  \pause
  \[
    x^2 + 4y^2
    \pause
    = 0
    \pause
    \neq 8
  \]
  \pause
  Não resolve o sistema

  \pause
  \vspace{\baselineskip}
  Portanto \(y\neq 0\)
\end{resolution}

%------------------------------------------------------------------------------%
\begin{resolution}{Casos 2 e 3}
\begin{columns}
\column{0.5\linewidth}
  \[
    y \neq 0
    \qquad
    \lambda = 2
  \]
  \[
    \pause
    x
    = \lambda y
    \pause
    = 2 y
  \]
  \[
    \pause
    x^2 = 4y^2
  \]
  \pause
  \[
    x^2 + 4y^2
    \pause
    = 4y^2 + 4y^2
    \pause
    = 8y^2
    \pause
    = 8
  \]
  \[
    \pause
    y^2 = 1
  \]
  \pause
  Portanto \(y=\pm 1\)
  \pause
\column{0.5\linewidth}
  \[
    y \neq 0
    \qquad
    \lambda = -2
  \]
  \[
    \pause
    x
    = \lambda y
    \pause
    = -2 y
  \]
  \[
    \pause
    x^2 = 4y^2
  \]
  \pause
  \[
    x^2 + 4y^2
    \pause
    = 4y^2 + 4y^2
    \pause
    = 8y^2
    \pause
    = 8
  \]
  \[
    \pause
    y^2 = 1
  \]
  \pause
  Portanto \(y=\pm 1\)
\end{columns}
  \pause
  \vspace{\baselineskip}
  Temos então os pontos
  \(( 2,  1)\),
  \(( 2, -1)\),
  \((-2,  1)\) e
  \((-2, -1)\)
\end{resolution}

%------------------------------------------------------------------------------%
\begin{resolution}{Solução}
  \onslide<+->{Avaliando a função \(f(x, y)=xy\) nos pontos}
  \begin{align*}
    \onslide<+->{f(-2, -1)}
    \onslide<+->{& = (-2)(-1)}
    \onslide<+->{  =  2 \\}
    \onslide<+->{f(\hphantom{-}2, -1)}
    \onslide<+->{& = 2 (-1)}
    \onslide<+->{  = -2 \\}
    \onslide<+->{f(-2, \hphantom{-}1)}
    \onslide<+->{& = -2 \times 1}
    \onslide<+->{  = -2 \\}
    \onslide<+->{f(\hphantom{-}2, \hphantom{-}1)}
    \onslide<+->{& =  2 \times 1}
    \onslide<+->{  = 2}
  \end{align*}

  \onslide<+->{O valor máximo de $f$ é $\hpm 2$
    e ocorre nos pontos \((-2, -1)\) e \(( 2, 1)\)}

  \onslide<+->{O valor mínimo de $f$ é $-2$
    e ocorre nos pontos \((-2, 1)\) e \((2, -1)\)}
\end{resolution}

%------------------------------------------------------------------------------%

%------------------------------------------------------------------------------%
\addtocounter{exemplo}{1}
\begin{frame}{Exemplo \theexemplo}
  Encontre os valores máximo e mínimo da função
  \[
    f(x, y) = 3x + 4y
  \]
  na circunferência
  \[
    x^2 + y^2 = 1
  \]
\end{frame}

%------------------------------------------------------------------------------%
\begin{frame}{Exemplo \theexemplo}
  Queremos encontrar os valores extremos da função
  \[
    f(x, y) = 3x + 4y
  \]
  \pause
  sujeito a restrição
  \[
    g(x, y) = x^2 + y^2 - 1 = 0
  \]

  \pause
  Para isso precisamos encontrar os valores $x$, $y$, $z$ e $\lambda$ tais que
  \[
    \nabla f = \lambda \nabla g
    \qquad\qquad
    g(x, y) = 0
  \]
\end{frame}

%------------------------------------------------------------------------------%
\begin{frame}{Exemplo \theexemplo}
\begin{columns}
\column{0.4\linewidth}
  Gradiente de $f$
  \[
    \pause
    \D{f}{x}
    \pause
    = \D{}{x}\left(3x + 4y\right)
    \pause
    = 3
  \]
  \[
    \pause
    \D{f}{y}
    \pause
    = \D{}{y}\left(3x + 4y\right)
    \pause
    = 4
  \]
  \pause
  \[
    {\color<17>{blue}\nabla f = \vetor{3}{4}}
  \]
\column{0.6\linewidth}
  \pause
  Gradiente de $g$
  \[
    \pause
    \D{g}{x}
    \pause
    = \D{}{x}\left(x^2 + y^2 - 1\right)
    \pause
    = 2x
  \]
  \[
    \pause
    \D{g}{y}
    \pause
    = \D{}{y}\left(x^2 + y^2 - 1\right)
    \pause
    = 2y
  \]
  \pause
  \[
    {\color<17>{blue}\nabla g = \vetor{2x}{2y}}
  \]
\end{columns}
\end{frame}

%------------------------------------------------------------------------------%
\begin{frame}[t]{Exemplo \theexemplo}
\begin{columns}
\column{0.5\linewidth}
  A equação
  \[
    \nabla f = \lambda \nabla g
  \]
  \pause
  se torna
  \[
    \vetor{3}{4} = \lambda \vetor{2x}{2y}
  \]
  {\color<7>{blue}%
    \pause
    \[
    \begin{cases}
      x = \dfrac{3}{2\lambda} \\[4mm]
      y = \dfrac{2}{\lambda}
    \end{cases}
  \]}
\column{0.5\linewidth}
  \pause
  A equação
  \[
    g(x, y) = 0
  \]
  \pause
  se torna
  \[
    x^2 + y^2 - 1 = 0
  \]
  {\color<7>{blue}%
    \pause
  \[
    x^2 + y^2 = 1
  \]}
\end{columns}
\end{frame}

%------------------------------------------------------------------------------%
\begin{frame}{Exemplo \theexemplo}
\begin{columns}
\column{0.5\linewidth}
  \onslide<+->{Precisamos resolver o sistema
  \[
  \begin{cases}
    x = \dfrac{3}{2\lambda} \\[3mm]
    y = \dfrac{2}{\lambda} \\[3mm]
    x^2 + y^2 = 1
  \end{cases}
  \]}
\column{0.5\linewidth}
  \begin{align*}
    \onslide<+->{x^2 + y^2 & = 1 \\[2mm]}
    \onslide<+->{\left(\frac{3}{2\lambda}\right)^2 +
                 \left(\frac{2}{\lambda}\right)^2 & = 1  \\[2mm]}
    \onslide<+->{\frac{9}{4\lambda^2} +
                 \frac{4}{\lambda^2} & = 1  \\[2mm]}
    \onslide<+->{\frac{9+16}{4\lambda^2} & = 1  \\[2mm]}
    \onslide<+->{4\lambda^2 & = 25  \\[2mm]}
    \onslide<+->{\lambda & = \pm \frac{5}{2}}
  \end{align*}
\end{columns}
\end{frame}

%------------------------------------------------------------------------------%
\begin{frame}{Exemplo \theexemplo}
\begin{columns}
\column{0.5\linewidth}
  Quando \(\lambda=\sfrac{5}{2}\)
  \[
    \pause
    x
    = \frac{3}{2\lambda}
    \pause
    = \frac{3}{2\sfrac{5}{2}}
    \pause
    = \frac{3}{5}
  \]
  \[
    \pause
        y
    = \frac{2}{\lambda}
    \pause
    = \frac{2}{\sfrac{5}{2}}
    \pause
    = \frac{4}{5}
  \]
\column{0.5\linewidth}
  \pause
  Quando \(\lambda=-\sfrac{5}{2}\)
  \[
    \pause
    x
    = \frac{3}{2\lambda}
    \pause
    = -\frac{3}{2\sfrac{5}{2}}
    \pause
    = -\frac{3}{5}
  \]
  \[
    \pause
    y
    = \frac{2}{\lambda}
    \pause
    = -\frac{2}{\sfrac{5}{2}}
    \pause
    = -\frac{4}{5}
  \]
\end{columns}
\end{frame}

%------------------------------------------------------------------------------%
\begin{frame}{Exemplo \theexemplo}
  \onslide<+->{Avaliando a função \(f(x, y)=3x + 4y\) nos pontos}
  \begin{align*}
    \onslide<+->{f\left(           -\frac{3}{5},            -\frac{4}{5}\right) &= }
    \onslide<+->{3\left(-\frac{3}{5}\right) + 4\left(-\frac{4}{5}\right) = }
    \onslide<+->{\frac{-9-16}{5} = }
    \onslide<+->{-5 \\[5mm]}
    \onslide<+->{f\left(\hphantom{-}\frac{3}{5}, \hphantom{-}\frac{4}{5}\right) &= }
    \onslide<+->{3\frac{3}{5} + 4\frac{4}{5} = }
    \onslide<+->{\frac{9+16}{5} = }
    \onslide<+->{5}
  \end{align*}

  \onslide<+->{O valor máximo de $f$ é \hspace{1ex} 5 e ocorre no ponto
  \(\left(\frac{3}{5}, \frac{4}{5}\right)\)}

  \vspace{0.3\baselineskip}
  \onslide<+->{O valor mínimo de $f$ é $-5$ e ocorre no ponto
  \(\left(-\frac{3}{5}, -\frac{4}{5}\right)\)}
\end{frame}

%------------------------------------------------------------------------------%

%------------------------------------------------------------------------------%
\begin{question}
  Encontre o ponto da superfície \(z=xy+1\) que está mais próximo da origem
\end{question}

%------------------------------------------------------------------------------%
\begin{resolution}{Formulação do Problema}
  Queremos minimizar a distância até a origem
  \[
    d(x, y, z) = \sqrt{x^2 + y^2 + z^2}
  \]
  \pause
  Como a função raiz quadrada é crescente, podemos minimizar a função
  \[
    f(x, y, z) = x^2 + y^2 + z^2
  \]
  \pause
  sujeito a restrição
  \[
    g(x, y, z) = z - xy = 1
  \]
\end{resolution}

%------------------------------------------------------------------------------%
\begin{resolution}{Gradientes}
\begin{columns}
\column{0.5\linewidth}
  \onslide<+->{Gradiente de \(f(x, y, z) = x^2 + y^2 + z^2\)}
  \begin{align*}
    \onslide<+->{f_x(x, y, z) & = 2x \\}
    \onslide<+->{f_y(x, y, z) & = 2y \\}
    \onslide<+->{f_z(x, y, z) & = 2z}
  \end{align*}

  \onslide<+->{
  \[
    \nabla f = \vetortri{2x}{2y}{2z}
  \]}
\column{0.5\linewidth}
  \onslide<+->{Gradiente de \(g(x, y, z)=z-xy\)}
  \begin{align*}
    \onslide<+->{g_x(x, y, z) & = -y \\}
    \onslide<+->{g_y(x, y, z) & = -x \\}
    \onslide<+->{g_z(x, y, z) & = 1}
  \end{align*}
  \onslide<+->{
  \[
    \nabla g = \vetortri{-y}{-x}{1}
  \]}
\end{columns}
\end{resolution}

%------------------------------------------------------------------------------%
\begin{resolution}{Sistema}
  Aplicando os Multiplicadores de Lagrange, precisamos resolver o sistema
  \[
  \begin{cases}
    2x = -\lambda y \\
    2y = -\lambda x \\
    2z = \lambda  \\
    z-xy = 1
  \end{cases}
  \]
\end{resolution}

%------------------------------------------------------------------------------%
\begin{resolution}{Solução}
  Isolamos $x$ na primeira equação e substituímos na segunda
  \[
    \pause
    2x = -\lambda y
    \pause
    \hspace{54mm}
    x = \frac{-\lambda y}{2}
  \]
  \pause
  \[
    y
    \pause = \frac{-\lambda}{2} x
    \pause = \frac{-\lambda}{2} \; \frac{-\lambda y}{2}
    \pause = \frac{\lambda^2}{4} y
    \hspace{20mm}
    \pause 4y = \lambda^2 y
  \]

  \pause
  \vspace{\baselineskip}
  Temos dois casos \(y=0\) ou \(y\neq 0\)
\end{resolution}

%------------------------------------------------------------------------------%
\begin{resolution}{Caso 1}
  Caso \(y=0\)
  \[
    \pause
    x = \frac{-\lambda y}{2}
    \pause
    = 0
  \]

  \pause
  \(z-xy = 1\) se reduz a \(z=1\)

  \pause
  \vspace{\baselineskip}
  Obtemos o ponto \((x_1, y_1, z_1) = (0, 0, 1)\)
\end{resolution}

%------------------------------------------------------------------------------%
\begin{resolution}{Caso 2}
  Caso \(y\neq 0\)
  \[
    \pause
    4y = \lambda^2 y
  \]
  \[
    \pause
    \lambda^2 = 4
  \]
  \[
    \pause
    \lambda = \pm 2
  \]
\end{resolution}

%------------------------------------------------------------------------------%
\begin{resolution}{Caso 2}
  \begin{columns}[t]
  \column{0.5\linewidth}
    Se $y\neq0$ e $\lambda=2$
    \[
      \pause x
      \pause = \frac{-\lambda y}{2}
      \pause = \frac{-2y}{2}
      \pause = -y
    \]

    \[
      \pause 2z = \lambda
      \pause = 2
    \]
    \[
      \pause z = 1
    \]
  \column{0.5\linewidth}
    \pause
    Substituindo na quarta equação
    \pause
    \begin{align*}
      \onslide<+->{z-xy      & = 1 \\}
      \onslide<+->{1 - (-y)y & = 1 \\}
      \onslide<+->{y^2       & = 0 \\}
      \onslide<+->{y         & = 0   }
    \end{align*}
    \pause
    Contradição
  \end{columns}
\end{resolution}

%------------------------------------------------------------------------------%
\begin{resolution}{Caso 2}
  \begin{columns}[t]
  \column{0.5\linewidth}
    Se $y\neq 0$ e $\lambda=-2$
    \[
      \pause x
      \pause = \frac{-\lambda y}{2}
      \pause = \frac{-(-2)y}{2}
      \pause = y
    \]

    \[
      \pause 2z = \lambda = -2
    \]
    \[
      \pause z = -1
    \]
  \column{0.5\linewidth}
    \pause
    Substituindo na quarta equação
    \pause
    \begin{align*}
      \onslide<+->{z-xy     & = 1  \\}
      \onslide<+->{-1 - y^2 & = 1  \\}
      \onslide<+->{-y^2     & = 2  \\}
      \onslide<+->{y^2      & = -2}
    \end{align*}
  \end{columns}

  \vspace{\baselineskip}
  \onslide<+->{Não existe solução}
\end{resolution}

%------------------------------------------------------------------------------%
\begin{resolution}{Solução}
  O ponto mais próximo da origem é \((0, 0, 1)\)
\end{resolution}

%------------------------------------------------------------------------------%

%------------------------------------------------------------------------------%
\addtocounter{exemplo}{1}
\begin{frame}{Exemplo \theexemplo}
  Encontre o ponto do plano \(x+2y+3z=13\) mais próximo do ponto \((1, 1, 1)\)
\end{frame}

%------------------------------------------------------------------------------%
\begin{frame}{Exemplo \theexemplo}
  Queremos minimizar a distância
  \[
    d(x, y, z) = \sqrt{(x-1)^2 + (y-1)^2 + (z-1)^2}
  \]
  Como a função raiz quadrada é crescente podemos minimizar a função
  \[
    f(x, y, z) = (x-1)^2 + (y-1)^2 + (z-1)^2
  \]
  suas derivadas parciais são
  \[
    f_x(x, y, z) = 2(x-1)
    \qquad
    f_y(x, y, z) = 2(y-1)
    \qquad
    f_z(x, y, z) = 2(z-1)
  \]
\end{frame}

%------------------------------------------------------------------------------%
\begin{frame}{Exemplo \theexemplo}
  As derivadas parciais da função \(g(x, y, z)=x+2y+3z\) são
  \[
    g_x(x, y, z) = 1
    \qquad
    g_y(x, y, z) = 2
    \qquad
    g_z(x, y, z) = 3
  \]
\end{frame}

%------------------------------------------------------------------------------%
\begin{frame}{Exemplo \theexemplo}
  Aplicando os Multiplicadores de Lagrange, precisamos resolver o sistema
  \[
  \begin{cases}
    2(x-1) =  \lambda \\
    2(y-1) = 2\lambda \\
    2(z-1) = 3\lambda \\
    x + 2y + 3z = 13
  \end{cases}
  \]
\end{frame}

%------------------------------------------------------------------------------%
\begin{frame}{Exemplo \theexemplo}
  Das três primeiras equações obtemos
  \[
    2(x-1) = \lambda
    \quad\Rightarrow\quad
    x = \frac{\lambda}{2} + 1
  \]
  \[
    2(y-1) = 2\lambda
    \quad\Rightarrow\quad
    y = \lambda + 1
  \]
  \[
    2(z-1) = 3\lambda
    \quad\Rightarrow\quad
    z = \frac{3\lambda}{2} + 1
  \]
\end{frame}

%------------------------------------------------------------------------------%
\begin{frame}{Exemplo \theexemplo}
  Substituindo na última equação
  \begin{align*}
    x + 2y + 3z & = 13
    \\[2mm]
    \left(\frac{\lambda}{2} + 1\right)
    + 2 \left(\lambda + 1\right)
    + 3 \left(\frac{3\lambda}{2} + 1\right) & = 13
    \\[2mm]
    \frac{\lambda}{2} + 1
    + 2\lambda + 2
    + \frac{9\lambda}{2} + 3 & = 13
    \\[2mm]
    \lambda + 2
    + 4\lambda + 4
    + 9\lambda + 6 & = 13\times 2
    \\[2mm]
    14\lambda + 12 & = 26
    \\[2mm]
    14\lambda & = 14
    \\[2mm]
    \lambda & = 1
  \end{align*}
\end{frame}

%------------------------------------------------------------------------------%
\begin{frame}{Exemplo \theexemplo}
  Portanto
  \[
    x = \frac{\lambda}{2} + 1 = \frac{3}{2}
  \]
  \[
    y = \lambda + 1 = 2
  \]
  \[
    z = \frac{3\lambda}{2} + 1 = \frac{5}{2}
  \]
  O ponto mais próximo é o ponto \(\left(\frac{3}{2}, 2, \frac{5}{2}\right)\)
\end{frame}

%------------------------------------------------------------------------------%

%------------------------------------------------------------------------------%
\addtocounter{exemplo}{1}
\begin{frame}{Exemplo \theexemplo}
  Encontre o ponto da esfera \(x^2+y^2+z^2=4\) mais distante do ponto \((1, -1, 1)\)

  \textit{Dica: não é necessário avaliar completamente a distância para determinar o máximo}
\end{frame}

%------------------------------------------------------------------------------%
\begin{frame}{Exemplo \theexemplo}
  Queremos minimizar a distância
  \[
    d(x, y, z) = \sqrt{(x-1)^2 + (y+1)^2 + (z-1)^2}
  \]
  Como a função raiz quadrada é crescente podemos minimizar a função
  \[
    f(x, y, z) = (x-1)^2 + (y+1)^2 + (z-1)^2
  \]
  suas derivadas parciais são
  \[
    f_x(x, y, z) = 2(x-1)
    \qquad
    f_y(x, y, z) = 2(y+1)
    \qquad
    f_z(x, y, z) = 2(z-1)
  \]
\end{frame}

%------------------------------------------------------------------------------%
\begin{frame}{Exemplo \theexemplo}
  As derivadas parciais da função \(g(x, y, z)=x^2+y^2+z^2\) são
  \[
    g_x(x, y, z) = 2x
    \qquad
    g_y(x, y, z) = 2y
    \qquad
    g_z(x, y, z) = 2z
  \]
\end{frame}

%------------------------------------------------------------------------------%
\begin{frame}{Exemplo \theexemplo}
  Aplicando os Multiplicadores de Lagrange, precisamos resolver o sistema
  \[
  \begin{cases}
    2(x-1) = 2\lambda x \\
    2(y+1) = 2\lambda y \\
    2(z-1) = 2\lambda z \\
    x^2 + y^2 + z^2 = 4
  \end{cases}
  \]
\end{frame}

%------------------------------------------------------------------------------%
\begin{frame}{Exemplo \theexemplo}
  Das três primeiras equações obtemos
  \begin{align*}
    2(x-1)       & = 2\lambda x          & 2(y+1)       & = 2\lambda y          & 2(z-1)       & = 2\lambda z          \\
    x-1          & = \lambda x           & y+1          & = \lambda y           & z-1          & = \lambda z           \\
    x(1-\lambda) & = 1                   & y(1-\lambda) & = -1                  & z(1-\lambda) & = 1                   \\
    x            & = \frac{1}{1-\lambda} & y            & =\frac{-1}{1-\lambda} & z            & = \frac{1}{1-\lambda}
  \end{align*}
\end{frame}

%------------------------------------------------------------------------------%
\begin{frame}{Exemplo \theexemplo}
  Sabemos que \(1-\lambda\neq 0\) pois \(x(1-\lambda) = 1\).

  Substituindo na última equação
  \begin{align*}
    x^2 + y^2 + z^2 & = 4
    \\[2mm]
    \left(\frac{1}{1-\lambda}\right)^2 +
    \left(\frac{-1}{1-\lambda}\right)^2 +
    \left(\frac{1}{1-\lambda}\right)^2 & = 4
    \\[2mm]
    \frac{1 + 1 + 1}{\left(1-\lambda\right)^2} & = 4
    \\[2mm]
    \frac{1}{\left(1-\lambda\right)^2} & = \frac{4}{3}
    \\[2mm]
    \frac{1}{1-\lambda} & = \pm\sqrt{\frac{4}{3}}
    = \pm\frac{2}{\sqrt{3}}
  \end{align*}
\end{frame}

%------------------------------------------------------------------------------%
\begin{frame}{Exemplo \theexemplo}
  Para \(\frac{1}{1-\lambda} = \frac{2}{\sqrt{3}}\)
  \[
    x = \frac{2}{\sqrt{3}}
    \qquad
    y = \frac{-2}{\sqrt{3}}
    \qquad
    z = \frac{2}{\sqrt{3}}
  \]
\end{frame}

%------------------------------------------------------------------------------%
\begin{frame}{Exemplo \theexemplo}
  Para \(\frac{1}{1-\lambda} = -\frac{2}{\sqrt{3}}\)
  \[
    x = \frac{-2}{\sqrt{3}}
    \qquad
    y = \frac{2}{\sqrt{3}}
    \qquad
    z = \frac{-2}{\sqrt{3}}
  \]
\end{frame}

%------------------------------------------------------------------------------%
\begin{frame}{Exemplo \theexemplo}
  Temos dois candidatos
  \(\left( \frac{2}{\sqrt{3}}, \frac{-2}{\sqrt{3}}, \frac{2}{\sqrt{3}}\right)\) e
  \(\left( \frac{-2}{\sqrt{3}}, \frac{2}{\sqrt{3}}, \frac{-2}{\sqrt{3}}\right)\).
\end{frame}

%------------------------------------------------------------------------------%
\begin{frame}{Exemplo \theexemplo}
  Avaliando a função em cada ponto temos
  \begin{align*}
    f\left( \frac{2}{\sqrt{3}}, \frac{-2}{\sqrt{3}}, \frac{2}{\sqrt{3}}\right)
    & = \left(\frac{ 2}{\sqrt{3}}-1\right)^2
    + \left(\frac{-2}{\sqrt{3}}+1\right)^2
    + \left(\frac{ 2}{\sqrt{3}}-1\right)^2 \\[2mm]
    & = \left(\frac{2}{\sqrt{3}}-1\right)^2
    + \left(\frac{2}{\sqrt{3}}-1\right)^2
    + \left(\frac{2}{\sqrt{3}}-1\right)^2 \\[2mm]
    & = 3\left(\frac{2}{\sqrt{3}}-1\right)^2
  \end{align*}
\end{frame}

%------------------------------------------------------------------------------%
\begin{frame}{Exemplo \theexemplo}
  \begin{align*}
    f\left( \frac{-2}{\sqrt{3}}, \frac{2}{\sqrt{3}}, \frac{-2}{\sqrt{3}}\right)
    & = \left(\frac{-2}{\sqrt{3}}-1\right)^2
    + \left(\frac{ 2}{\sqrt{3}}+1\right)^2
    + \left(\frac{-2}{\sqrt{3}}-1\right)^2 \\[2mm]
    & = \left(\frac{2}{\sqrt{3}}+1\right)^2
    + \left(\frac{2}{\sqrt{3}}+1\right)^2
    + \left(\frac{2}{\sqrt{3}}+1\right)^2 \\[2mm]
    & = 3\left(\frac{2}{\sqrt{3}}+1\right)^2
  \end{align*}
\end{frame}

%------------------------------------------------------------------------------%
\begin{frame}{Exemplo \theexemplo}
  Portanto o ponto mais distante é
  \(\left( \frac{-2}{\sqrt{3}}, \frac{2}{\sqrt{3}}, \frac{-2}{\sqrt{3}}\right)\)
  \clearpage
\end{frame}

%------------------------------------------------------------------------------%


%------------------------------------------------------------------------------%
\section{Lista Mínima}

%------------------------------------------------------------------------------%
\begin{frame}{Lista Mínima}
  Cálculo Vol. 2 do Thomas 12\textsuperscript{a} ed. -- Seção 14.8
  \begin{enumerate}
    \item Estudar o texto da seção
    \item Resolver os exercícios: 3, 5, 7, 12, 16, 18, 24 e 30
  \end{enumerate}

  \vfill
  Atenção: A prova é baseada no livro, não nas apresentações
\end{frame}

%------------------------------------------------------------------------------%
\end{document}
%------------------------------------------------------------------------------%
